% ============================================================
% 第五讲 5.8：财富效应、商品分类与综合例题
% ============================================================
\subsection{财富效应与 Engel 曲线}

Slutsky 方程把价格效应拆成替代效应与财富效应，两块的符号由不同力量决定。本节先按\textbf{财富效应}的方向给商品分类，再按\textbf{自身价格效应}的方向分类，并用 Slutsky 方程把两套分类焊起来。

\begin{definition}[正常品与劣等品]
固定 $(p,w)$，设 $x_i(p,w)$ 可微。若 $\partial x_i/\partial w>0$（财富增加消费增加），称商品 $i$ 在 $(p,w)$ 处为\textbf{正常品}（normal good）；若 $\partial x_i/\partial w<0$，称为\textbf{劣等品}（inferior good）。若对所有 $w$ 都是正常品，称全局正常品。
\end{definition}

劣等品并非“质量差”，而是\textbf{随财富增长被排挤出消费束}的商品：便宜的主食、低档交通方式等——穷时多买，富时少买。正常与劣等是\emph{局部}性质：同一商品可能在低收入区间正常、高收入区间劣等。

财富 $w$ 变动（价格不动）时，最优束 $x(p,w)$ 在商品空间中描出的轨迹称为\textbf{收入--消费曲线}（income--consumption curve, ICC）；把其中第 $i$ 个分量对 $w$ 作图，即得第 $i$ 种商品的 \textbf{Engel 曲线} $w\mapsto x_i(p,w)$。

\begin{figure}[H]
\centering
\begin{tikzpicture}[scale=.9]
  \draw[axisline] (0,0)--(6.1,0) node[below]{$x_1$};
  \draw[axisline] (0,0)--(0,4.75) node[left]{$x_2$};
  % 三条预算线 x+y=w, w=2.4,3.6,4.8（p=(1,1)）
  \draw[curveC] (0,2.4)--(2.4,0);
  \draw[thick] (0,3.6)--(3.6,0);
  \draw[curveB] (0,4.8)--(4.8,0);
  \node[figlabel] at (1.05,2.62) {$w_1$};
  \node[figlabel] at (2.25,3.82) {$w_2$};
  \node[figlabel] at (3.55,5.02) {$w_3$};
  % 三条无差异曲线（u=x^{1/3}y^{2/3} 型，切点 (0.8,1.6),(1.2,2.4),(1.6,3.2)）
  \draw[curveA,domain=0.55:3.1,smooth] plot (\x,{1.431/sqrt(\x)});
  \draw[curveA,domain=0.75:4.3,smooth] plot (\x,{2.629/sqrt(\x)});
  \draw[curveA,domain=0.95:5.4,smooth] plot (\x,{4.048/sqrt(\x)});
  % 收入-消费曲线：斜率 2 的射线
  \draw[thick,-{Stealth}] (0.15,0.3)--(2.0,4.0);
  \node[figlabel,above left] at (1.95,3.95) {ICC};
  \node[nodewhite] at (0.8,1.6) {};
  \node[figlabel,right] at (0.87,1.55) {$x^{(1)}$};
  \node[nodewhite] at (1.2,2.4) {};
  \node[figlabel,right] at (1.27,2.35) {$x^{(2)}$};
  \node[nodewhite] at (1.6,3.2) {};
  \node[figlabel,right] at (1.67,3.15) {$x^{(3)}$};
\end{tikzpicture}
\caption{财富增加：收入--消费曲线与 Engel 曲线}
\figexplan{价格固定为 $p=(1,1)$，财富由 $w_1$ 增至 $w_3$，预算线（两条虚点线夹一条实线）平行外移，最优束沿切点 $x^{(1)}\to x^{(2)}\to x^{(3)}$ 移动，其轨迹即收入--消费曲线 ICC（带箭头粗线）。此处两种商品都是正常品：ICC 单调向右上。把 $x_1$ 分量对 $w$ 作图即第一种商品的 Engel 曲线；对位似偏好（本图效用为 $x_1^{1/3}x_2^{2/3}$），ICC 是过原点的射线，Engel 曲线是过原点的直线。}
\label{fig:engel}
\end{figure}

\begin{remark}[Engel 曲线的形状与位似偏好]
位似偏好（边际替代率只依赖消费比例）下 $x(p,w)=w\cdot\phi(p)$，Engel 曲线是过原点的直线；其中 Cobb--Douglas 给出 $x_i=\alpha_i w/p_i$，斜率即支出份额除以价格。恩格尔本人观察到食品支出份额随财富上升而下降（Engel 定律），即食品的财富弹性小于 $1$——效用函数层面常用食品份额内嵌的 PIGLOG/QUAIS 形式来拟合这一特征。
\end{remark}

\subsection{自身价格效应：普通品与 Giffen 品}

\begin{definition}[普通品与 Giffen 品]
若 $\partial x_i/\partial p_i<0$（自身涨价、需求减少），称商品 $i$ 为\textbf{普通品}（ordinary good）；若 $\partial x_i/\partial p_i>0$（自身涨价、需求反而增加），称为\textbf{Giffen 品}。
\end{definition}

两套分类通过 Slutsky 方程严格相连。

\begin{proposition}[正常品必为普通品]\label{prop:normal-ordinary}
若商品 $i$ 在 $(p,w)$ 处正常（$\partial x_i/\partial w>0$）且 $x_i(p,w)>0$，则 $\partial x_i/\partial p_i<0$，即 $i$ 是普通品。
\end{proposition}
\begin{proof}
由 Slutsky 方程（定理~\ref{thm:slutsky-eq}）取 $i=j$：
\[
 \frac{\partial x_i}{\partial p_i}
 =\frac{\partial h_i}{\partial p_i}-x_i\,\frac{\partial x_i}{\partial w}.
\]
右边第一项由补偿需求定律（推论~\ref{cor:diag}）非正；第二项中 $x_i>0$、$\partial x_i/\partial w>0$，故 $-x_i\,\partial x_i/\partial w<0$。两项之和严格为负。
\end{proof}

\begin{corollary}[Giffen 品必为劣等品]
若商品 $i$ 是 Giffen 品（$\partial x_i/\partial p_i>0$）且 $x_i>0$，则 $\partial x_i/\partial w<0$。
\end{corollary}
\begin{proof}
命题~\ref{prop:normal-ordinary} 的逆否命题：若 $\partial x_i/\partial w\ge0$ 且 $x_i>0$，则 $\partial x_i/\partial p_i<0$，不可能是 Giffen 品。
\end{proof}

因此 Giffen 品必须同时满足苛刻的三个条件：(i) 劣等品；(ii) 在预算中占\textbf{很大}份额（财富效应 $\propto x_i$ 才够大）；(iii) 替代余地小（补偿效应弱）。三者叠加使“涨价使你更穷 $\Rightarrow$ 更多依赖这种廉价主食”的机制压倒替代机制。

\begin{remark}[爱尔兰土豆故事的评述]
教科书的经典叙事是：$19$ 世纪爱尔兰饥荒中土豆价格上涨，贫穷家庭实际购买力锐减，被迫\emph{减少}肉与面包、\emph{增加}土豆消费，土豆遂成 Giffen 品。评述有三点。其一，该故事主要是 Marshall 的转述，Giffen 本人的著述中并无土豆的明确记载，作为史实并不可靠。其二，机制本身在逻辑上成立且已被实验支持：Jensen 与 Miller（2008）在中国湖南与甘肃对大米/面食的补贴实验中观察到了与 Giffen 行为一致的证据——对象恰是“收入不高、该主食占预算比重大、无更便宜的替代品”的家庭。其三，Giffen 行为是\textbf{局部}现象：同一商品只在贫富区间的一段内、且只在特定价格范围内表现如此；它违背的是马歇尔需求定律，而不是补偿需求定律（后者由偏好凸性保证永不可能被违反）。这再次说明把“需求定律”精确化为 Slutsky 分解的价值：可违反的从来只是财富效应那一项。
\end{remark}

\begin{figure}[H]
\centering
\begin{tikzpicture}[scale=.88]
  \draw[axisline] (0,0)--(7.2,0) node[below]{$x_1$};
  \draw[axisline] (0,0)--(0,4.75) node[left]{$x_2$};
  % 原预算线 (0,4)-(4,0)
  \draw[thick] (0,4)--(4,0);
  \node[figlabel] at (3.6,4.3) {原预算线};
  % 新预算线 (0,4)-(8,0)，画到 (7.4,0.3)
  \draw[thick,dashed] (0,4)--(7.4,0.3);
  \node[figlabel] at (6.0,4.42) {新预算线};
  % 补偿预算线 (0,2.83)-(5.66,0)
  \draw[thick,dotted] (0,2.83)--(5.66,0);
  \node[figlabel] at (5.1,3.3) {补偿预算线};
  % 无差异曲线 u1: xy=4
  \draw[curveA,domain=1.0:5.3,smooth] plot (\x,{4/\x});
  \node[figlabel] at (4.7,1.3) {$u_1$};
  % 无差异曲线 u2：在 (1.5,3.25) 与新预算线相切（分段抛物线拼接，相切于 x=1.5）
  \draw[curveA,domain=1.15:1.5,smooth] plot (\x,{4-0.5*\x+2*(\x-1.5)*(\x-1.5)});
  \draw[curveA,domain=1.5:4.8,smooth] plot (\x,{4-0.5*\x+0.15*(\x-1.5)*(\x-1.5)});
  \node[figlabel] at (1.02,4.35) {$u_2$};
  % 三点
  \node[nodeblue] at (2,2) {};
  \node[figlabel,below] at (2,1.85) {$x$};
  \node[nodewhite] at (2.83,1.41) {};
  \node[figlabel,right] at (2.93,1.36) {$h$};
  \node[nodeblue] at (1.5,3.25) {};
  \node[figlabel,above left] at (1.42,3.32) {$x'$};
  % 分解箭头
  \draw[arrblue] (2.3,1.95) to[bend left=32] (2.62,1.5);
  \node[figlabel] at (3.3,1.05) {替代效应（$x_1\uparrow$）};
  \draw[arrblue] (2.6,1.6) to[bend left=45] (1.75,2.9);
  \node[figlabel] at (3.35,2.35) {财富效应（$x_1\downarrow\downarrow$，占优）};
\end{tikzpicture}
\caption{Giffen 品：财富效应压倒替代效应}
\figexplan{商品 1 劣等且占预算大头。$p_1$ 下降使预算线外旋（实线 $\to$ 虚线）。补偿点 $h$ 在原无差异曲线 $u_1$ 上：替代效应使 $x_1$ 由 $2$ 增至 $2.83$。但真实财富效应把消费者推向更高的无差异曲线 $u_2$ 上的 $x'$，其横坐标 $1.5$ \emph{小于}原点 $x$ 的横坐标 $2$——降价反而少买，即 Giffen 行为。注意 $u_2$ 在左侧急剧抬升：劣等品的财富效应方向与替代效应相反，且数值上占优。}
\label{fig:giffen}
\end{figure}

\subsection{交叉价格效应：总替代与 Hicksian 替代}

自身价格之外的另一组分类看\textbf{交叉}价格效应 $\partial x_i/\partial p_j$（$i\ne j$），并区分“是否已补偿”。

\begin{definition}[总替代/总互补与 Hicksian 替代/互补]
商品 $j$ 是商品 $i$ 的\textbf{总替代品}（gross substitute），若 $\partial x_i/\partial p_j>0$；\textbf{总互补品}（gross complement），若 $\partial x_i/\partial p_j<0$（均指未补偿的 Marshallian 需求）。相应地，若补偿需求满足 $\partial h_i/\partial p_j>0$，称 $j$ 是 $i$ 的 \textbf{Hicksian 替代品}；若 $\partial h_i/\partial p_j<0$，称 \textbf{Hicksian 互补品}。
\end{definition}

两套概念的关键差别在\textbf{对称性}：

\begin{remark}[总概念不对称，补偿概念对称]
总替代不必对称：$\partial x_i/\partial p_j$ 与 $\partial x_j/\partial p_i$ 可以一正一负。原因在 Slutsky 方程的财富项：
\[
 \frac{\partial x_i}{\partial p_j}
 =\frac{\partial h_i}{\partial p_j}-x_j\frac{\partial x_i}{\partial w},\qquad
 \frac{\partial x_j}{\partial p_i}
 =\frac{\partial h_j}{\partial p_i}-x_i\frac{\partial x_j}{\partial w},
\]
两个财富项一般不相等（取决于各自的 $x$ 与 $\partial x/\partial w$）。而替代项由 Slutsky 矩阵的\textbf{对称性}（定理~\ref{thm:slutsky-props}）保证 $\partial h_i/\partial p_j=\partial h_j/\partial p_i$：补偿意义上“$i$ 的替代品”与“$j$ 的替代品”是同一关系，像“邻居”一样是相互的。这使 Hicksian 分类在理论上更干净；总分类则更贴近市场上可观察的未补偿反应，二者各有用途，但不应混用。
\end{remark}

\begin{example}[Cobb--Douglas：总分类为中性，补偿分类为替代]
$u=x_1^{\alpha}x_2^{1-\alpha}$ 给出 $x_1=\alpha w/p_1$：$\partial x_1/\partial p_2=0$，两商品\textbf{既非总替代也非总互补}；但补偿意义上 $\partial h_1/\partial p_2=\alpha(1-\alpha)w/(p_1p_2)>0$（见下节例~\ref{ex:cd-slutsky}），是严格的 Hicksian 替代品。同一对商品在两套分类下结论不同，直观原因：CD 下两种商品的支出份额固定，$p_2$ 变动不改变花在商品 $1$ 上的钱（总效应为零），但固定效用后相对价格的替代机制仍然存在。
\end{example}

\begin{proposition}[两商品世界必为 Hicksian 替代品]\label{prop:two-goods}
设 $n=2$，偏好凸、$h$ 可微且 $D_p h$ 负半定、行和为零。则 $\partial h_1/\partial p_2\ge0$ 且 $\partial h_2/\partial p_1\ge0$；若 $\partial h_1/\partial p_1<0$（自身补偿效应严格为负），则两者严格为正。
\end{proposition}
\begin{proof}
对第 $1$ 行写出行和为零（Euler）：$h_{11}p_1+h_{12}p_2=0$，故
\[
 h_{12}=-\frac{p_1}{p_2}\,h_{11}\;\ge\;0,
\]
末一不等号用了对角元非正（推论~\ref{cor:diag}）。再由对称性 $h_{21}=h_{12}\ge0$。严格性直接由 $h_{11}<0$ 与上式传递。
\end{proof}

直觉：只有两种商品时，保持效用不变、商品 $1$ 涨价，钱只能省在商品 $2$ 上（或反过来）——补偿世界里的“敌人”只有对方，必然互为替代。互补品（如咖啡与糖）至少需要三种商品：商品 $1$、商品 $2$，以及涨价后把两者一起挤出的“第三方”。

\subsection{Offer 曲线}

\begin{definition}[价格--提供曲线]
固定财富 $w$ 与 $p_{-i}$，让 $p_i$ 变动，最优束 $x(p,w)$ 在商品空间中描出的轨迹称为\textbf{价格--提供曲线}（price offer curve）。
\end{definition}

提供曲线是 Slutsky 分解的“原始底片”：它把价格变化引起的\emph{全部}总效应记录成一条路径，每一点的切向即 $\nabla_p x$ 沿该方向的列。提供曲线向哪边弯，直接显示商品是普通品还是 Giffen 品、总替代还是总互补。

\begin{figure}[H]
\centering
\begin{tikzpicture}[scale=.88]
  \draw[axisline] (0,0)--(7.5,0) node[below]{$x_1$};
  \draw[axisline] (0,0)--(0,4.35) node[left]{$x_2$};
  % 三条预算线均过 (0,4)：p1=2,1,0.5
  \draw[curveC] (0,4)--(2,0);
  \node[figlabel] at (1.05,4.25) {$p_1=2$};
  \draw[thick] (0,4)--(4,0);
  \node[figlabel] at (3.05,4.25) {$p_1=1$};
  \draw[curveB] (0,4)--(7.4,0.3);
  \node[figlabel] at (5.9,4.35) {$p_1=\tfrac12$};
  % 提供曲线 x2=4/(x1+1)
  \draw[curveA,domain=0.25:7.25,smooth] plot (\x,{4/(\x+1)});
  \node[figlabel] at (5.9,1.0) {offer 曲线};
  % 一条无差异曲线过 (3,1)：x1=3-ln(x2)
  \draw[curveB,domain=0.55:3.0,smooth,variable=\t] plot ({3-ln(\t)},{\t});
  \node[figlabel] at (1.6,2.5) {$u=3$};
  % 切点
  \node[nodewhite] at (1,2) {};
  \node[figlabel,above left] at (0.9,2.1) {$x^{(2)}$};
  \node[nodewhite] at (3,1) {};
  \node[figlabel,above left] at (2.9,1.1) {$x^{(1)}$};
  \node[nodewhite] at (7,0.5) {};
  \node[figlabel,left] at (6.9,0.55) {$x^{(1/2)}$};
\end{tikzpicture}
\caption{价格--提供曲线（拟线性效用 $u=x_1+\ln x_2$）}
\figexplan{财富固定为 $w=4$、$p_2=1$，$p_1$ 由 $2$ 降到 $1$ 再降到 $1/2$，预算线绕纵截距 $(0,4)$ 外旋（点线 $\to$ 实线 $\to$ 虚线）。切点 $x^{(2)}\to x^{(1)}\to x^{(1/2)}$ 的轨迹即提供曲线（实粗线），此处为双曲线 $x_2=4/(x_1+1)$。虚线无差异曲线 $u=3$ 与中间预算线相切于 $x^{(1)}=(3,1)$，可检验切点性质。提供曲线向右下弯：$p_1$ 下降时 $x_1$ 沿曲线增加、$x_2$ 减少，商品 $1$ 是普通品，且商品 $2$ 是商品 $1$ 的总互补品（$\partial x_2/\partial p_1<0$ 对称地看即 $\partial x_2$ 随 $p_1$ 下降而下降）。}
\label{fig:offer}
\end{figure}

\subsection{综合例题}

\begin{example}[拟线性效用的完整求解与 Slutsky 验证]\label{ex:quasilinear}
设 $u(x_1,x_2)=x_1+\ln x_2$，$p\gg0$。最大化 $x_1+\ln x_2$ 于 $p_1x_1+p_2x_2=w$。
\textbf{需求。}内点一阶条件：$1=\lambda p_1$，$1/x_2=\lambda p_2$，故
\[
 x_2=\frac{p_1}{p_2},\qquad
 x_1=\frac{w-p_1}{p_2}\cdot\frac{p_2}{p_1}=\frac{w}{p_1}-1 .
\]
内点解要求 $w>p_1$（否则 $x_1=0$，全部财富买 $x_2=w/p_2$——非负约束在 $w\le p_1$ 时绑定，是角点解）。
\textbf{间接效用。}
\[
 v(p,w)=\frac{w}{p_1}-1+\ln\frac{p_1}{p_2}.
\]
Roy identity 验证：$v_w=1/p_1=\lambda$；$v_{p_1}=-w/p_1^2+1/p_1$，$v_{p_2}=-1/p_2$，故
\[
 -\frac{v_{p_1}}{v_w}=\frac{w}{p_1}-1=x_1,\qquad
 -\frac{v_{p_2}}{v_w}=\frac{p_1}{p_2}=x_2 .\checkmark
\]
\textbf{支出函数。}由对偶恒等式，令 $v(p,w)=\bar u$ 反解 $w$：
\[
 e(p,\bar u)=p_1\Bigl(\bar u+1-\ln\frac{p_1}{p_2}\Bigr).
\]
\textbf{Hicksian 需求。}Shephard 引理：
\[
 h_1=\frac{\partial e}{\partial p_1}=\bar u-\ln\frac{p_1}{p_2},\qquad
 h_2=\frac{\partial e}{\partial p_2}=\frac{p_1}{p_2}.\checkmark
\]
（$h_2$ 与 Marshallian $x_2$ 相同：商品 $2$ 无财富效应，补偿与否无差别。）
\textbf{Slutsky 验证。}先算 Marshallian 导数：$\partial x_1/\partial p_1=-w/p_1^2$，$\partial x_1/\partial w=1/p_1$，$x_1=w/p_1-1$。补偿导数：由 $h_1=\bar u-\ln p_1+\ln p_2$ 得（在 $\bar u=v$ 处）$\partial h_1/\partial p_1=-1/p_1$。代入 Slutsky 方程右端：
\[
 \frac{\partial h_1}{\partial p_1}-x_1\frac{\partial x_1}{\partial w}
 =-\frac{1}{p_1}-\Bigl(\frac{w}{p_1}-1\Bigr)\frac{1}{p_1}
 =-\frac{w}{p_1^2}
 =\frac{\partial x_1}{\partial p_1} .\checkmark
\]
交叉项：$\partial x_1/\partial p_2=0$；$\partial h_1/\partial p_2=1/p_2$，$x_2=p_1/p_2$，$\partial x_1/\partial w=1/p_1$，故右端 $=1/p_2-(p_1/p_2)(1/p_1)=0$ $\checkmark$。对称性：$\partial h_1/\partial p_2=1/p_2=\partial h_2/\partial p_1$ $\checkmark$。矩阵性质：由 $h_2=p_1/p_2$ 补齐第二行，得
\[
 S=\begin{pmatrix}
 -1/p_1 & 1/p_2\\
 1/p_2 & -p_1/p_2^2
 \end{pmatrix},
\]
行列式 $=(1/p_1)(p_1/p_2^2)-1/p_2^2=0$、对角元均为负，故特征值为 $0$ 与负数，负半定 $\checkmark$；两行与 $p$ 内积均为零，$Sp=0$ $\checkmark$。
\end{example}

\begin{example}[Cobb--Douglas 的 Slutsky 矩阵]\label{ex:cd-slutsky}
设 $u=x_1^{\alpha}x_2^{1-\alpha}$，$0<\alpha<1$。已知
\[
 x_1=\frac{\alpha w}{p_1},\qquad x_2=\frac{(1-\alpha)w}{p_2},\qquad
 e(p,\bar u)=\frac{\bar u\,p_1^{\alpha}p_2^{1-\alpha}}{\alpha^{\alpha}(1-\alpha)^{1-\alpha}} .
\]
\textbf{补偿需求。}$h_1=\partial e/\partial p_1=\alpha e/p_1$，$h_2=(1-\alpha)e/p_2$。
\textbf{Slutsky 矩阵。}在 $\bar u=v(p,w)$（此时 $e=w$）处：
\[
 \frac{\partial h_1}{\partial p_1}
 =-\alpha(1-\alpha)\frac{e}{p_1^2}\Big|_{e=w}
 =-\frac{\alpha(1-\alpha)w}{p_1^2},\qquad
 \frac{\partial h_1}{\partial p_2}
 =\frac{\alpha(1-\alpha)w}{p_1p_2},
\]
由对称结构（$\alpha\leftrightarrow1-\alpha$、$p_1\leftrightarrow p_2$）：
\[
 \frac{\partial h_2}{\partial p_1}=\frac{\alpha(1-\alpha)w}{p_1p_2},\qquad
 \frac{\partial h_2}{\partial p_2}=-\frac{\alpha(1-\alpha)w}{p_2^2}.
\]
于是
\[
 S(p,w)=-\alpha(1-\alpha)w
 \begin{pmatrix}
 1/p_1^2 & -1/(p_1p_2)\\
 -1/(p_1p_2) & 1/p_2^2
 \end{pmatrix}
 =-\alpha(1-\alpha)w\;q\,q^{\top},\qquad
 q=\Bigl(\frac1{p_1},-\frac1{p_2}\Bigr)^{\top}.
\]
\textbf{三大性质验证。}
对称性：$S_{12}=S_{21}$，显然。负半定：对任意 $z\in\R^2$，
\[
 z^{\top}Sz=-\alpha(1-\alpha)w\,(q^{\top}z)^2\le0,
\]
且 $S$ 秩为 $1$，特征值为 $0$ 与 $-\alpha(1-\alpha)w\|q\|^2<0$。行和为零：$q^{\top}p=p_1/p_1-p_2/p_2=0$，故 $Sp=-\alpha(1-\alpha)w\,q\,(q^{\top}p)=0$。
\textbf{Slutsky 方程验证。}自身价格：
\[
 \frac{\partial h_1}{\partial p_1}-x_1\frac{\partial x_1}{\partial w}
 =-\frac{\alpha(1-\alpha)w}{p_1^2}-\frac{\alpha w}{p_1}\cdot\frac{\alpha}{p_1}
 =-\frac{\alpha w}{p_1^2}=\frac{\partial x_1}{\partial p_1}\;\checkmark
\]
交叉价格：$\partial x_1/\partial p_2=0$，而
\[
 \frac{\partial h_1}{\partial p_2}-x_2\frac{\partial x_1}{\partial w}
 =\frac{\alpha(1-\alpha)w}{p_1p_2}-\frac{(1-\alpha)w}{p_2}\cdot\frac{\alpha}{p_1}=0\;\checkmark.
\]
注意结论：CD 的 $S$ 负半定但不满秩（一行乘比例即另一行），体现两商品世界中补偿替代关系的秩约束；而 Marshallian 交叉效应为零，再次对照出总分类与补偿分类的差别。
\end{example}
